% =============================================================================
% 013 / 068
% Status: raw
% =============================================================================
% Notes:
%
% Source question:
% 請問，歷史上曾經有從上古或中古漢語分化出的齊語、楚語和蜀語嗎？
% =============================================================================

\chapter[請問，歷史上曾經有從上古或中古漢語分化出的齊語、楚語和蜀語嗎？]{請問，歷史上曾經有從上古或中古漢語分化出的齊語、楚語和蜀語嗎？}
\qaanswer{
這個我不清楚，不過我知道以下3個關於四川地區語言變化脈絡的知識點：

1.在北宋中葉有士大夫的文獻中記載，蜀語(梁益方言)類閩。也就是說當時的蜀語和當時閩地的語言有很多近似的地方。但北宋時的蜀語和現在的四川話差異很大，整個面貌完全不同。這涉及到宋代到清初四川地區的數次人口置換。

2.現在的四川話顯然是清中期人口佔絕對優勢的兩湖人、贛人填四川時覆蓋過去的西南官話和少數碎片化了的老湘語、贛語。它們與之前傳入四川的江淮軍話，殘餘的老蜀語一起混合穩定後顯了現在西南官話四川區片的狀態。

3.在清以前的明代整個西南地區的各殖民點大城就已經被作為明初軍事征服者的朱元璋的江淮部隊「軍話」(江淮官話)和徽州移民講的徽語(高地吳語的一種)覆蓋過一次了。

2、3的情況導致大的西南官話區的形成，現在四川、貴州、雲南、湖北、湖南的語言互通率很高。而且越是晚期被漢地風俗所同化的少數民族地區西南官話使用率反而越高。比如湘黔交界的定居苗族不使用老湘語反而使用西南官話。雲南、貴州很多土司地區的情況也類似，越晚被同化的地區講的官話反而越「地道」。
}

